\documentclass[11pt,a4paper]{article}
\usepackage[utf8]{inputenc}
\usepackage[T1]{fontenc}
\usepackage[margin=2cm]{geometry}
\usepackage{booktabs}
\usepackage{enumitem}
\usepackage{xcolor}
\usepackage{eurosym}
\usepackage{titlesec}
\usepackage{fancyhdr}
\usepackage{microtype}

\definecolor{artiso}{HTML}{4a4280}
\definecolor{muted}{HTML}{6B7280}

\titleformat{\section}{\large\bfseries\color{artiso}}{}{0em}{}[\vspace{-4pt}\color{artiso}\rule{\textwidth}{0.5pt}]
\titleformat{\subsection}{\normalsize\bfseries}{}{0em}{}

\pagestyle{fancy}
\fancyhf{}
\renewcommand{\headrulewidth}{0pt}
\fancyfoot[C]{\small\color{muted}ARTISO — Speaker Scripts — April 9, 2026 — CONFIDENTIAL}
\fancyfoot[R]{\small\color{muted}\thepage}

\setlength{\parindent}{0pt}
\setlength{\parskip}{6pt}

\begin{document}

\begin{center}
{\LARGE\bfseries\color{artiso} ARTISO — Speaker Scripts}\\[6pt]
{\large Board of Directors Presentation — April 9, 2026}\\[4pt]
{\normalsize 10 minutes + 5 min Q\&A}\\[12pt]
\begin{tabular}{llrl}
\toprule
\textbf{Speaker} & \textbf{Slides} & \textbf{Time} & \textbf{Focus} \\
\midrule
Leonor & 1--2 & 50s & Opening + Board framing \\
Mart\'{i} & 3--4 & 2:25 & P\&L + Revenue Engine \\
Chris & 5--7 & 2:20 & Unit Econ + Cash + Financing \\
Roxy & 8--10 & 2:40 & Cap Table + Valuation + Scenarios \\
Tom\'{a}s & 11--12 & 1:20 & Risks + Board Ask \\
\midrule
& & \textbf{9:35} & \textit{25s buffer for transitions} \\
\bottomrule
\end{tabular}
\end{center}

\vspace{8pt}

% ============================================================
\section{Leonor --- Slides 1--2 (50s)}
% ============================================================

\subsection{Slide 1: ARTISO at a Glance (\texttildelow25s)}

Good morning, we welcome you to ARTISO --- a Barcelona-based AI platform that helps fashion brands design and launch collections faster, cutting development time by 70\% and campaign costs by up to 90\%, built by former Mango executives. Seven paying customers including Deichmann, over \euro{}60K in pilot revenue, and 200+ brands on the waitlist.

\subsection{Slide 2: Board Decisions (\texttildelow25s)}

Today we need the board to decide three things: how much capital ARTISO needs to reach Series~A from strength, what financing structure protects founder control, and whether our \euro{}27M pre-money valuation holds up.

Our recommendation: raise \euro{}3M through a blended stack at 4.5\% cost of capital --- 60\% cheaper than pure equity. Mart\'{i}, the numbers.

% ============================================================
\section{Mart\'{i} --- Slides 3--4 (2:25)}
% ============================================================

\subsection{Slide 3: 5-Year P\&L Summary (\texttildelow80s)}

Thank you Leonor. Here is our 5-year P\&L, built entirely bottom-up from sales capacity and pricing assumptions --- not from top-down TAM capture.

Revenue grows from \euro{}1.2M in FY2026 to \euro{}36.5M by FY2030. Gross margins hold steady around 82\% throughout, reflecting our SaaS-plus-credits model where the main cost driver is cloud compute.

On the cost side, we invest heavily in years one and two --- total OpEx reaches \euro{}4.5M in FY2027 as we scale the team from 4 to 23. EBITDA starts at negative \euro{}755K, narrows through FY2027 and 2028, and we hit EBITDA breakeven in Q4~2028. By FY2030, EBITDA reaches \euro{}8.2M at a 22.5\% margin.

The installed base grows from 35 brands in year one to 360 by FY2030. End-of-year ARR reaches \euro{}35.7M. And importantly, cumulative pre-profitability losses of approximately \euro{}2.1M are shielded by tax loss carry-forward --- so the company pays zero taxes until FY2029.

One thing to note: FY2026 starts in April when the pre-seed closes. Q1~2026 is pre-investment with zero revenue and zero operating costs.

\subsection{Slide 4: Revenue Engine (\texttildelow65s)}

Let me show you how we built the revenue model. This is the engine underneath the P\&L.

We start with founder selling time --- 0.85 FTEs --- and scale to 7 hired account executives by year five. Each AE closes approximately 2 deals per month, validated against Salesloft benchmarks. Sales cycles range from 1 month for SME, 3 months for Large, to 6 months for Enterprise.

The funnel splits across three tiers. SME brands at \euro{}17K annual contract value, Large at \euro{}67K, and Enterprise at \euro{}155K. Revenue comes from four streams: one-off setup fees, recurring subscriptions, usage-based credits --- which is the primary growth driver --- and professional services capped below 10\%.

On the right you can see how ACV expands after landing. Customers start with one module and a base subscription, then cross-sell into additional modules and increase credit usage over time. The ARR milestones track this: \euro{}100K ARR by Q3~2026, the critical \euro{}1M ARR gate by Q3~2027 which triggers Series~A readiness, and \euro{}35.7M by Q4~2030. Chris, over to you for unit economics.

% ============================================================
\section{Chris --- Slides 5--7 (2:20)}
% ============================================================

\subsection{Slide 5: Unit Economics (\texttildelow45s)}

Our unit economics confirm this model is sustainable. Blended CAC starts at \euro{}17K with founder-led sales and rises to \euro{}74K as we scale into enterprise --- that's a realistic trajectory, not fantasy metrics.

LTV-to-CAC lands at 5.3x, well above the 3x threshold. Payback is 11 months. Burn multiple trends to zero as we approach profitability. And the Rule of 40 reaches 64\% by FY2030 --- 41.5\% revenue growth plus 22.5\% EBITDA margin.

\subsection{Slide 6: Funding Need \& Timing (\texttildelow45s)}

The cash waterfall tells the story. After the pre-seed closes, cash depletes as we invest in growth. The critical point: cash dips to negative \euro{}379K in Q4~2028 --- the pre-Series~A gap. That's why timing matters.

Series~A prep must begin by Q4~2027 --- 15 months of lead time. And a venture debt window opens in H2~2028 near EBITDA breakeven as an additional safety net. The \euro{}3M pre-seed funds 24 months of runway.

\subsection{Slide 7: Blended Financing Stack (\texttildelow50s)}

Now the financing strategy --- this is our key recommendation. We are not proposing a pure equity raise.

The lead VC contributes \euro{}1.5 to \euro{}2M at 5--7\% dilution. Angels and SAFEs bridge \euro{}300 to \euro{}500K. Then we layer non-dilutive capital on top: ENISA at \euro{}300K, CDTI NEOTEC at \euro{}250 to \euro{}400K, and Catalan grants at \euro{}200 to \euro{}300K. Total: approximately \euro{}3M at only 8--12\% dilution.

The comparison is clear: 4.5\% blended cost of capital versus 11\% for pure equity. That's 60\% cheaper --- and saves founders 2 to 5 percentage points of ownership that they keep through exit. Roxy will show you the cap table impact.

% ============================================================
\section{Roxy --- Slides 8--10 (2:40)}
% ============================================================

\subsection{Slide 8: Cap Table \& Dilution (\texttildelow50s)}

Pre-round, each founder holds 33.3\%. After the pre-seed, founders retain 26.7\% each --- 53.3\% combined --- preserving majority control. ESOP is set at 10\%, and pre-seed investors receive 10\%.

Post-Series~A, founders dilute to 24.2\% each, with Series~A investors taking 9.4\%. Price per share goes from \euro{}2.40 at pre-seed to \euro{}7.72 at Series~A --- a 3.2x step-up. We use the VC-standard pre-money ESOP carve-out, noting that the pro-rata method would be fairer to founders.

\subsection{Slide 9: Valuation --- 5-Step VC Method (\texttildelow70s)}

We justify \euro{}27M pre-money through the 5-step VC method.

Step one: FY2030 ARR of \euro{}35.7M at a 10x exit multiple gives \euro{}357M enterprise value. Step two: at a 60\% IRR hurdle --- appropriate for pre-seed risk --- the required investor value at exit is \euro{}19.7M. Step three: that requires 5.5\% ownership. Steps four and five give us an implied pre-money of \euro{}51.4M.

So at \euro{}27M, investors enter at a 47\% discount to the VC-implied valuation --- pricing in execution risk.

The triangulation on the right confirms the range. Forward revenue gives \euro{}21.4M at the low end, revenue and ARR multiples give \euro{}41 to \euro{}49M, and comparables anchor at \euro{}27M. The range is \euro{}21.4 to \euro{}51.4M --- our \euro{}27M sits comfortably within it.

Pre-seed investor IRR: 85.7\%, with an 11.9x return multiple. That far exceeds the 60\% hurdle.

\subsection{Slide 10: Scenarios \& Sensitivity (\texttildelow40s)}

Three scenarios stress-test the plan. Base case: \euro{}36.5M revenue by FY2030, Series~A in Q4~2027. Upside: faster enterprise adoption accelerates everything. Downside: we activate contingency at the \euro{}500K cash trigger --- defer 3 hires, cut S\&M 30\%.

The sensitivity tornado shows two dominant levers: Enterprise NRR and lead growth rate. This is a volume-driven model. Everything else is secondary. Tom\'{a}s will cover the risks.

% ============================================================
\section{Tom\'{a}s --- Slides 11--12 (1:20)}
% ============================================================

\subsection{Slide 11: Risks \& Mitigation (\texttildelow50s)}

Five key risks. Two are critical: customer concentration --- top 3 logos are over 60\% of year-one ARR --- mitigated by diversifying across tiers and clustering through MODACC. And technology moat erosion if Adobe or OpenAI enter vertical fashion AI --- mitigated by deepening the proprietary fashion ontology and PLM integrations.

Two are high: extended enterprise sales cycles, mitigated by a parallel SME pipeline; and cash runway timing, mitigated by starting the Series~A process in Q4~2027 with a \euro{}500K internal trigger. The contingency plan is ready: cut burn from \euro{}120K to \euro{}80K per month, extending runway by 5 months. Venture debt at H2~2028 provides a further bridge.

\subsection{Slide 12: Board Ask (\texttildelow30s)}

Three resolutions. One: approve the \euro{}3M pre-seed at \euro{}27M pre-money with the blended structure. Two: launch the VC process and file ENISA and CDTI in parallel, starting immediately. Three: begin Series~A preparation in Q4~2027.

The discipline is not just raising this round --- it is ensuring the company reaches the next one with evidence, optionality, and leverage. Thank you. We welcome your questions.

\end{document}
